\documentclass[11pt,a4paper]{article}
\usepackage[T1]{fontenc}
\usepackage[utf8]{inputenc}
\usepackage{lmodern}
\usepackage{amsmath}
\usepackage{booktabs}
\usepackage{graphicx}
\usepackage{xcolor}
\usepackage{geometry}
\usepackage{microtype}
\usepackage{hyperref}
\usepackage{array}
\geometry{margin=24mm}
\hypersetup{colorlinks=true,linkcolor=blue!50!black,urlcolor=blue!50!black,
  pdftitle={Parallel Chess Search with MPI, OpenMP, and GPU Offload},
  pdfauthor={Jia Siqi and Jia Zeyue}}
\setlength{\parskip}{0.35em}
\setlength{\parindent}{1.2em}
\newcommand{\code}[1]{\texttt{#1}}

\title{Parallel Chess Search with MPI, OpenMP, and GPU Offload}
\author{Jia Siqi (A0333724H) \and Jia Zeyue (A0326500U)}
\date{DSA5208 Project 3 --- draft, 27 October 2026}

\begin{document}
\maketitle

\begin{abstract}
We study fixed-depth chess search as an irregular parallel workload. A serial
negamax implementation with alpha-beta pruning is compared with OpenMP root
splitting, deeper MPI+OpenMP task frontiers, a seeded-root variant that
retains more pruning, and OpenMP GPU offload. The same legal-move library and
simple material evaluation are used throughout. On Vanda, a depth-8,
split-depth-3 task search visited 929.9 million nodes. A one-worker run took
198.1 seconds; a ten-node, 720-core feasibility run reported 1.59 seconds of
internal search time. The latter is a single trial pending a repeated run.
The serial alpha-beta baseline visited only 25.6 million nodes and took a
median of 5.34 seconds. These contrasting baselines reveal that increasing
parallelism can lose so much pruning that hardware speedup over the serial
algorithm remains modest. An A40 GPU executed the offload kernel correctly,
but the measured depth-3 and depth-4 searches were far slower than serial.
\end{abstract}

\section{Problem and algorithm}

\subsection{What a chess search computes}

A chess position records the pieces on 64 squares, the player to move,
castling rights, an en-passant square, and counters needed for the draw
rules. At each position the player has a set of legal moves. Applying a move
changes the position and hands the turn to the opponent. Repeating this
process yields a game tree: a node is a position and an edge is a legal move.
The branching factor varies greatly with the position. Moves may also create
checkmate, stalemate, or a draw. We use Forsyth--Edwards notation (FEN) to
specify reproducible starting positions.

Our program does not aim to play competitive chess. It searches to a fixed
ply depth $d$, where one ply is one move by one player. At a nonterminal
depth-zero leaf, it scores the material balance using pawn = 100, knight =
320, bishop = 330, rook = 500, and queen = 900. The king is not assigned a
material value because it cannot be captured. Positive scores favor the side
to move. Checkmate has a large negative score adjusted by distance to mate;
stalemate and supported draw conditions score zero. A one-pawn advantage is
thus roughly 100 score units. The leaf evaluation ignores king safety,
positional value, quiescence, and opening knowledge. A score or selected move
should therefore be read as a result of this experimental search, not as a
claim about the best move in chess.

\subsection{Minimax, negamax, and alpha-beta pruning}

Chess is a two-player zero-sum game: a gain for the player to move is a loss
for the opponent. A minimax search assumes that both players choose the best
available continuation. Negamax writes the same recurrence with one sign
change at each turn:
\begin{equation}
  V(s,d)=\begin{cases}
    E(s), & d=0\text{ or }s\text{ is terminal},\\
    \displaystyle\max_{m\in M(s)} -V(s_m,d-1), & \text{otherwise},
  \end{cases}
\end{equation}
where $s_m$ is the position after legal move $m$. Terminal outcomes are
handled before the material evaluation. The root returns a score and the
first legal move attaining the maximum. A simple exhaustive minimax backend
serves as an independent search-control reference at shallow depth: it uses
the same legal moves and evaluation but disables pruning.

Alpha-beta pruning avoids searching continuations that cannot affect the
choice above them \cite{knuth1975}. The search carries a lower bound
$\alpha$ and an upper bound $\beta$. A child is searched with the negated
window $[-\beta,-\alpha]$; if $\alpha\geq\beta$, later siblings cannot
improve the relevant ancestor and are skipped. The returned value may be a
bound when a cutoff occurs. Move order matters: a strong move discovered
early can make many later branches cheap, while a poor first move leaves a
wide window. Our baseline uses the legal-move order supplied by the library,
without a transposition table or learned ordering. This makes the
dependency between sibling searches visible in the parallel experiments.

\section{Implementation and design decisions}

\subsection{Serial reference and third-party boundary}

The project implements search, evaluation, task construction, MPI/OpenMP
scheduling, GPU tree traversal, and counters in C++17. It vendors the
unmodified, MIT-licensed single-header \code{Disservin/chess-library}
revision \code{2a608d58f44776909bb6e2c09505e97ba1451319}
\cite{chesslib}. That library supplies board representation, FEN parsing,
legal move generation, move application and reversal, draw helpers, and UCI
move formatting; it does not supply our search algorithm. The instructor
approved this credited external dependency. The header and license are in
\code{third\_party/chess-library/}, with provenance in
\code{THIRD\_PARTY.md}.

\subsection{Independent root and deeper-frontier tasks}

The simplest CPU parallel method gives each legal root move an independent
full-window search. An OpenMP dynamic loop schedules the root moves among
threads; for MPI, moves are assigned to ranks by their index modulo the rank
count. Each task owns a board copy and private search counters. Afterward,
MPI combines candidate scores with \code{MPI\_Allreduce} and sums counters.
The search time reported for an MPI job is the maximum rank elapsed time,
including the local frontier construction, search, and reductions, but not
process launch or compilation.

A root may have only tens of moves, so root splitting cannot fill hundreds
of cores. We also materialize the game tree to split depth two or three.
Each frontier board becomes a task and performs an independent full-window
negamax search for the remaining depth. At the initial position, split depth
three produced 9,323 frontier nodes, enough tasks to distribute across 720
requested workers. The internal frontier nodes are combined bottom-up to
recover the exact minimax root result. This exposes more parallelism at a
cost: sibling tasks do not share alpha-beta bounds, so they can visit many
nodes that serial alpha-beta would prune. Each MPI rank currently builds the
same frontier. The assignment of tasks to ranks is static; OpenMP schedules
the rank's tasks dynamically.

\subsection{Seeded-root alternative}

To retain more pruning, a second root method first searches the initial
legal move with a full window. Its score $a$ is broadcast to the MPI ranks.
The remaining root moves use $a$ as a fixed bound. A task that fails to beat
$a$ may return an inexact bound, but it cannot change the root maximum; a
task that beats it supplies the needed value. This variant is implemented as
\code{omp-pv} and \code{mpi-pv}. It exposes only the remaining root moves,
and its static modulo rank assignment can be badly unbalanced. We measure
the seed time, seed node count, and per-rank parallel work to examine this
tradeoff.

\subsection{OpenMP target feasibility study}

Move generation is branch-heavy and uses a complex CPU board type. For the
GPU experiment, the CPU first expands the legal tree to the chosen shallow
depth and packs nodes and leaf evaluations into flat arrays. An OpenMP
\code{target teams distribute parallel for} region gives independent root
branches to the device \cite{openmp52}. Each device worker runs an iterative
stack-based alpha-beta traversal of its branch. This is a deliberately
bounded offload study: CPU tree expansion, device transfer, and limited root
parallelism remain part of the measured GPU path. We require
\code{OMP\_TARGET\_OFFLOAD=MANDATORY} and check
\code{omp\_is\_initial\_device()} to reject silent host fallback.

\section{Experimental method}

Runs used NUS Vanda PBS allocations, GCC 14.2.0 and Open MPI 5.0.7 for
CPU jobs, and NVHPC 24.7 with an A40 device for GPU jobs. The CPU script
requests one MPI rank per node and sets the indicated number of OpenMP
threads per rank. We initially discovered that Open MPI's default policy
bound a four-thread rank to one core. A controlled depth-7 rerun on the
same node fell from 6.23 to 2.43 seconds when \code{--bind-to none} was used
\cite{openmpi507}. A separate PBS affinity probe showed that the unbound
rank remained inside the four-CPU PBS cgroup. All corrected scaling runs use
this launcher setting. PBS \code{place=scatter} and logged hostnames verify
distinct nodes for multi-node jobs.

The main strong-scaling workload is the initial-position FEN, depth eight,
split depth three. It visits exactly 929,868,807 nodes in every measured
task-decomposition run. We repeat each 72-, 144-, and 288-core allocation
three times inside one PBS job and use the median internal search time.
The one-worker decomposition run and the 720-core feasibility run currently
have one trial each. A separate three-trial 720-core job is queued and will
replace the pilot point when complete. Source SHA-256, PBS job IDs, nodes,
toolchain libraries, scores, and raw output are preserved in
\code{experiments/raw/}. The independent serial alpha-beta baseline uses the
same initial position and depth but performs less work because it prunes
across siblings. We report both baselines rather than interpreting them as
the same scaling reference.
The exact main input is
\begin{quote}\small\ttfamily
rnbqkbnr/pppppppp/8/8/8/8/PPPPPPPP/RNBQKBNR w KQkq - 0 1
\end{quote}

\section{Results and interpretation}

\subsection{Strong scaling of the fixed task search}

For $p$ allocated CPU cores, we define speedup against a one-worker run of
the \emph{same task-decomposition algorithm} as
$S(p)=T_{\mathrm{task}}(1)/T_{\mathrm{task}}(p)$ and efficiency
$E(p)=S(p)/p$. Table~\ref{tab:scaling} and Figure~\ref{fig:scaling}
show internal search time, excluding MPI startup. The 720-core row is
provisional because it has one trial.

\begin{table}[ht]
\centering
\small
\caption{Initial position, depth 8, split depth 3. Search work is
929,868,807 nodes for every row. The one- and 720-core rows are single trials.}
\label{tab:scaling}
\begin{tabular}{@{}rrrrr@{}}
\toprule
Cores & Hosts & Time (s) & $S(p)$ & $E(p)$ \\
\midrule
1   & 1  & 198.129 & 1.00   & 100.0\% \\
72  & 1  & 9.094   & 21.79  & 30.3\% \\
144 & 2  & 3.048   & 65.01  & 45.1\% \\
288 & 4  & 1.945   & 101.86 & 35.4\% \\
720 & 10 & 1.593   & 124.36 & 17.3\% \\
\bottomrule
\end{tabular}
\end{table}

\begin{figure}[ht]
\centering
\includegraphics[width=\linewidth]{figures/scaling.pdf}
\caption{Vanda internal search time and efficiency. The open 720-core
marker is one trial, pending a three-trial result. The dashed time curve is
ideal scaling from the one-worker task search; the dotted line is the serial
alpha-beta median.}
\label{fig:scaling}
\end{figure}

The 72-to-144-core median ratio is 2.98 and the 72-to-288-core ratio is
4.68, both above ideal ratios. We do not attribute this to a particular
mechanism: allocation-specific node performance, cache effects, and the
static rank distribution are plausible contributors and were not isolated
in this series. Efficiency relative to the one-worker task search is only
17.3\% at the single 720-core point. The 720-core PBS job consumed 53
seconds of walltime although its internal search reported 1.593 seconds;
process launch, setup, and teardown matter for a complete end-to-end cost.

The serial alpha-beta search visited 25,604,258 nodes and took 5.33762,
5.34012, and 5.34224 seconds (median 5.34012). The task decomposition
visits about 36.3 times as many nodes. Thus the one-time 720-core search is
only 3.35 times faster than serial alpha-beta, despite using 720 allocated
cores. This is the central algorithmic cost of parallelizing branches
without their shared bounds.

\subsection{Root bounds versus exposed parallelism}

On the initial position at depth eight, the seeded-root method visited the
same 25,604,258 nodes as serial alpha-beta. Its medians were 3.010, 2.853,
and 4.205 seconds on one, two, and four four-thread nodes. Different node
types were assigned to these allocations, so the times are exploratory.
The instrumented four-node rerun on the same four CN hosts as the earlier
four-node job spent only 0.347 seconds on the seed move; afterward rank 0
searched 14.14 million nodes for 3.459 seconds, while the other ranks
searched 2.64--3.88 million nodes for 0.333--0.444 seconds. Static rank
assignment, rather than the seed search alone, dominated that critical path.
The comparison demonstrates that retaining pruning can reduce total work
while too few or poorly distributed tasks limit scaling.

A second, move-rich Kiwipete FEN at depth six checks whether the initial
position's result generalizes. Both serial and two-node seeded-root runs
selected \code{e2a6} with score $-90$. Their respective node counts were
19.23 and 26.89 million and median times 5.359 and 4.670 seconds. The
seeded-root method did more work here because its fixed initial bound was
weaker than the serial search's evolving root bound. Its two ranks received
6.90 and 19.39 million parallel nodes, another example of imbalance.
The exact Kiwipete input is recorded in \code{experiments/pv/README.md}
and the corresponding raw PBS logs.

\subsection{GPU execution and cost}

Both A40 jobs compiled and returned \code{ran\_on\_device=1}; GPU and
serial scores and selected moves agreed at depths three and four. Table
\ref{tab:gpu} separates the CPU tree build plus target path total from
time measured around the target regions. It does not separate device
launch, transfer, and kernel execution.

\begin{table}[ht]
\centering
\small
\caption{Initial-position A40 offload pilots, one trial each.}
\label{tab:gpu}
\begin{tabular}{@{}rrrrr@{}}
\toprule
Depth & CPU-built tree nodes & GPU total (s) & Target (s) & Serial (s) \\
\midrule
3 & 9,322   & 0.246888 & 0.244381 & 0.000125 \\
4 & 206,603 & 0.270310 & 0.238880 & 0.000237 \\
\bottomrule
\end{tabular}
\end{table}

The combined target-region time stayed near 0.24 seconds as the tree grew
more than twentyfold. The GPU total was roughly 1,974 and 1,138 times the
serial time at depths three and four. These small workloads and the
CPU-built full tree do not suit the GPU design; a larger GPU-native chess
representation and a different task strategy would be needed for a useful
performance comparison. We stop the offload study at this measured
limitation rather than claiming a GPU speedup.

\section{Validation, reproducibility, and limits}

The local smoke suite compares serial alpha-beta with exhaustive minimax,
OpenMP root, seeded-root OpenMP, deeper OpenMP frontiers, and the GPU kernel
on its host fallback path at depths one through three on seven FENs.
These include checkmate, stalemate, promotion, and en-passant positions.
All 21 position/depth comparisons passed. On Vanda, MPI root/frontier and
seeded-root runs matched the serial score on the initial position, the
seeded-root MPI search also matched serial on Kiwipete, and both A40 runs
matched serial at their measured depths. Equal-score moves may legitimately
differ; in the recorded comparisons the reported move also agreed. The
selected move is always drawn from the library's legal-move list.

These checks validate search-control agreement under the shared move
generator and evaluation; they are not an independent proof of full chess
rules or of deep tactical strength. The source, vendored library and license,
PBS scripts, raw logs, summary CSV, and figure source are shipped with this
report. Run \code{make test} locally, or submit the documented PBS scripts
after loading the recorded compiler modules. Logs identify source hashes and
job allocations. GPU compilation needs NVHPC plus the GCC 13.2 toolchain;
an initial failed compiler job is retained in the raw logs along with the
successful reruns.

\section{AI assistance and conclusion}

OpenAI Codex assisted with design discussion, C++ implementation, MPI and
OpenMP variants, GPU offload code, PBS scripts, SSH experiment execution,
diagnosis, result summaries, documentation, and drafting this report and
figure. The students chose the chess-search topic, supplied the assignment
and instructor guidance, requested cluster experiments, and must review the
code, measurements, interpretations, and final submission. We preserve raw
Vanda outputs so reported measurements can be checked against actual jobs.

The study illustrates a tension specific to tree search. Splitting more
deeply creates enough tasks for hundreds of cores, but destroys useful
alpha-beta dependencies and inflates total work. Searching a promising
root move first preserves pruning, yet exposes few tasks and suffers from
static load imbalance. GPU offload is technically possible but not useful
for this CPU-built, branch-heavy shallow tree. The strongest future
improvement would be a load-aware or dynamic MPI scheduler together with a
parallel search scheme that shares safe alpha-beta bounds. The queued
three-trial 720-core result will be incorporated before final submission.

\begin{thebibliography}{9}
\bibitem{knuth1975}
D. E. Knuth and R. W. Moore, ``An Analysis of Alpha-Beta Pruning,''
\emph{Artificial Intelligence}, vol. 6, no. 4, pp. 293--326, 1975.
\url{https://doi.org/10.1016/0004-3702(75)90019-3}.
\bibitem{chesslib}
Disservin, \emph{chess-library}, MIT-licensed C++ chess library, commit
\code{2a608d58f44776909bb6e2c09505e97ba1451319}.
\url{https://github.com/Disservin/chess-library}.
\bibitem{openmpi507}
Open MPI, \emph{mpirun / mpiexec}, version 5.0.7 documentation.
\url{https://docs.open-mpi.org/en/v5.0.7/man-openmpi/man1/mpirun.1.html}.
\bibitem{openmp52}
OpenMP Architecture Review Board, \emph{OpenMP API Specification}, version
5.2, 2021. \url{https://www.openmp.org/wp-content/uploads/OpenMP-API-Specification-5-2.pdf}.
\end{thebibliography}

\end{document}
